\documentclass[11pt,a4paper]{article}
\usepackage[T1]{fontenc}
\usepackage[utf8]{inputenc}
\usepackage[english]{babel}
\usepackage{amsmath,amssymb}
\usepackage[margin=2cm]{geometry}
\setlength{\parindent}{0pt}
\setlength{\emergencystretch}{3em}
\allowdisplaybreaks

\title{About \texttt{false\_equipotence.tex}}
\author{}
\date{}
\begin{document}
\maketitle
This note explains what the test data in \texttt{false\_equipotence.tex} is used for. That file is input for the engine in \texttt{main.c}; it is read by the test \texttt{test\_set\_theory.c} in this folder.

\begin{enumerate}
\item \textbf{Purpose.} Anti-regression for paragraph 17w of \texttt{parser.c} and 28q of \texttt{main.c}. The file holds 19 statements about equipotence that must not be proved, each a small change of a true one: items 1 to 4 of \texttt{infinite\_sets.tex} with a hypothesis dropped ($S \times T \sim S \times S$ without $S \sim T$, $\mathcal{P}(S) \sim \mathcal{P}(T)$ without it, $T_{1}^{S_{1}} \sim T_{2}^{S_{2}}$ without $T_{1} \sim T_{2}$, a product without $S_{1} \sim S_{2}$), a set of parts against a set ($\mathcal{P}(S) \sim T$, Cantor's $S \sim \mathcal{P}(S)$), base and exponent exchanged ($S^{V} \sim V^{T}$, $V^{S} \sim T^{W}$), a factor added or dropped ($S \times T \sim S$, $\mathcal{P}(S) \times V \sim \mathcal{P}(T)$, $S \times S \sim T$, $S \sim S \times S$, true only for infinite sets), a set unrelated to the hypotheses ($S \sim W$), negations and disjunctions that do not follow. Item 15, $\lnot (S \sim T) \implies \lnot (\mathcal{P}(S) \sim \mathcal{P}(T))$, holds for finite sets but is not a theorem: for infinite sets it cannot be proved (there are models of set theory with $2^{\aleph_{0}} = 2^{\aleph_{1}}$).

\item \textbf{What is being tested.} That none is proved. Every one but item 15 is refuted by a finite model, where $A \sim B$ is computed as ``the same number of elements'' (\texttt{model\_equipotent()}), with a certificate that \texttt{checker.c} recomputes (11j). Item 15 must stay \texttt{NON DIMOSTRATO}: the grammar of 28q never turns a bijection between sets of parts into one between the sets, and no rule of the engine counts elements. If the engine used the number of elements as a rule, or the grammar built a map that is not a bijection, or the checker accepted a definition of a wrong form, one of these would be at risk.

\item \textbf{How to run it.} From the top folder of the repository:
\begin{verbatim}
cc -std=c11 -O2 -Wall -Wextra -o /tmp/test_set_theory \
    tests/set_theory/test_set_theory.c && /tmp/test_set_theory
\end{verbatim}

\item \textbf{Expected result.} One \texttt{PASS} line for the file: none of the statements is proved; and two for the certificates of its counterexamples.
\end{enumerate}
\end{document}
